% 分类目录：ml
\documentclass[12pt, a4paper]{article}
\usepackage[margin=2.5cm]{geometry}
\usepackage{ctex}
\usepackage{amsmath, amssymb}
\usepackage{listings}
\usepackage{xcolor}
\usepackage{tabularx}
\usepackage{hyperref}
\usepackage{tikz}
\usetikzlibrary{shapes, arrows.meta, positioning}

\tikzset{
    treenode/.style={rectangle, draw, thick, minimum width=2.5cm, minimum height=0.7cm, align=center, rounded corners=2pt, fill=blue!8},
    leafnode/.style={rectangle, draw, thick, minimum width=1.8cm, minimum height=0.7cm, align=center, rounded corners=2pt, fill=green!15},
    treearrow/.style={-Stealth, thick}
}

\lstset{
    language=Python,
    basicstyle=\ttfamily\small,
    keywordstyle=\color{blue},
    commentstyle=\color{green!50!black},
    stringstyle=\color{red},
    showstringspaces=false,
    breaklines=true,
    numbers=left,
    numberstyle=\tiny\color{gray},
    frame=single
}

\title{知识点：决策树 (Decision Tree)}
\author{}
\date{}

\begin{document}
\maketitle

\section{核心定义}
\textbf{中文定义}：决策树是一种基于树形结构的监督学习算法。它通过递归地选择最优特征对数据集进行划分，从根节点到叶节点形成一系列\textbf{if-then}规则，最终实现分类或回归。

\textbf{英文定义}：A Decision Tree is a supervised learning algorithm that recursively partitions the feature space by selecting the most informative feature at each internal node, forming a tree of if-then rules for classification or regression.

\textbf{核心优势}：决策树是少数真正\textbf{可解释}的模型之一——人类可以直接阅读和理解其决策逻辑，在金融风控、医疗诊断等需要可解释性的领域不可替代。

\section{核心原理：如何选择最优划分特征？}
决策树构建的关键问题是：在每个节点上，选择哪个特征、以什么阈值进行划分？不同算法给出了不同的度量标准。

\subsection{信息增益 (Information Gain) — ID3}
\textbf{信息熵}衡量数据的不确定性：
$$ H(D) = -\sum_{k=1}^K p_k \log_2 p_k $$
\textbf{变量含义}：
\begin{itemize}
    \item $D$：当前数据集。
    \item $K$：类别总数。
    \item $p_k$：第 $k$ 类样本在 $D$ 中的占比。
    \item $H(D) \in [0, \log_2 K]$：熵越大，数据越"混乱"。
\end{itemize}
按特征 $A$ 划分后的\textbf{条件熵}：
$$ H(D|A) = \sum_{v=1}^V \frac{|D^v|}{|D|} H(D^v) $$
其中 $D^v$ 是特征 $A$ 取值为 $v$ 的子集，$V$ 是 $A$ 的取值数。

\textbf{信息增益}：
$$ \text{Gain}(D, A) = H(D) - H(D|A) $$
选择信息增益最大的特征进行划分。

\textbf{ID3 算法的缺陷}：信息增益偏向取值数目多的特征。例如，若存在一个"ID"特征（每个样本唯一），它的信息增益最大，但毫无泛化能力。

\subsection{信息增益率 (Gain Ratio) — C4.5}
C4.5 引入\textbf{分裂信息量}对信息增益进行惩罚：
$$ \text{SplitInfo}(D, A) = -\sum_{v=1}^V \frac{|D^v|}{|D|} \log_2 \frac{|D^v|}{|D|} $$
$$ \text{GainRatio}(D, A) = \frac{\text{Gain}(D, A)}{\text{SplitInfo}(D, A)} $$
分裂信息量 $\text{SplitInfo}$ 在特征取值多时会变大，从而自动惩罚"取值过多"的特征。

\subsection{基尼指数 (Gini Index) — CART}
基尼指数衡量从数据集中随机抽取两个样本、类别不一致的概率：
$$ \text{Gini}(D) = 1 - \sum_{k=1}^K p_k^2 $$
按特征 $A$ 划分后：
$$ \text{Gini}(D, A) = \sum_{v=1}^V \frac{|D^v|}{|D|} \text{Gini}(D^v) $$
选择使 $\text{Gini}(D, A)$ 最小的特征和划分点。CART 只进行\textbf{二叉划分}（将特征取值分为两组），每次产生恰好两个子节点。

\section{三大经典算法对比}
\begin{table}[h]
\centering
\begin{tabularx}{\textwidth}{|l|l|l|X|}
\hline
\textbf{算法} & \textbf{划分标准} & \textbf{树结构} & \textbf{核心特点} \\
\hline
ID3 & 信息增益 & 多叉树 & 只能处理离散特征，偏好多值特征。 \\
\hline
C4.5 & 信息增益率 & 多叉树 & 支持连续特征，修正了 ID3 的偏好问题。 \\
\hline
CART & 基尼指数 & 二叉树 & 既可分类也可回归，sklearn 默认实现。 \\
\hline
\end{tabularx}
\end{table}

\section{决策树示意图}
\begin{figure}[h]
\centering
\resizebox{0.75\textwidth}{!}{
\begin{tikzpicture}[node distance=1.5cm and 1.5cm]
    \node [treenode] (root) {天气 = ?};
    \node [leafnode, below left=of root, xshift=-1cm] (sunny) {晴天 $\to$ 不去};
    \node [treenode, below=of root] (cloudy) {温度 $> 30$?};
    \node [leafnode, below right=of root, xshift=1cm] (rainy) {雨天 $\to$ 不去};
    
    \node [leafnode, below left=of cloudy] (hot) {热 $\to$ 不去};
    \node [leafnode, below right=of cloudy] (cool) {凉爽 $\to$ 去};
    
    \draw [treearrow] (root) -- node[left, font=\small] {晴} (sunny);
    \draw [treearrow] (root) -- node[right, font=\small, xshift=-3mm] {阴} (cloudy);
    \draw [treearrow] (root) -- node[right, font=\small] {雨} (rainy);
    \draw [treearrow] (cloudy) -- node[left, font=\small] {是} (hot);
    \draw [treearrow] (cloudy) -- node[right, font=\small] {否} (cool);
\end{tikzpicture}}
\caption{决策树示意图：根据天气和温度决定是否出行}
\end{figure}

\section{剪枝 (Pruning)}
决策树如果不加限制地生长，会完美拟合训练数据（每个叶子节点只包含一个样本），导致严重\textbf{过拟合}。剪枝是决策树最重要的正则化手段。

\subsection{预剪枝 (Pre-pruning)}
在构建树的过程中提前停止分裂。常用策略：
\begin{itemize}
    \item 设定最大深度 (\texttt{max\_depth})。
    \item 设定叶子节点最小样本数 (\texttt{min\_samples\_leaf})。
    \item 设定分裂所需最小样本数 (\texttt{min\_samples\_split})。
    \item 设定最小信息增益阈值。
\end{itemize}
\textbf{优点}：训练速度快。\textbf{缺点}：可能过早停止分裂，导致欠拟合。

\subsection{后剪枝 (Post-pruning)}
先构建完整的决策树，再自底向上地评估每个子树：如果去掉子树后在验证集上的性能不下降（或仅微弱下降），则将该子树替换为叶子节点。
\textbf{优点}：通常能获得更好的泛化性能。\textbf{缺点}：需要额外的验证集，计算代价较高。

\section{CART 回归树}
CART 不仅能分类，还能回归。回归树在每个叶子节点输出该区域所有训练样本目标值的\textbf{均值}，划分标准改为最小化\textbf{均方误差 (MSE)}：
$$ \min_{A, s} \left[ \sum_{x_i \in D_1(A,s)} (y_i - c_1)^2 + \sum_{x_i \in D_2(A,s)} (y_i - c_2)^2 \right] $$
其中 $c_1, c_2$ 分别是左右子树的均值。

\section{集成学习扩展}
单棵决策树的能力有限，但通过集成方法可以显著提升性能：
\begin{itemize}
    \item \textbf{Bagging (Random Forest)}：多棵树并行训练，投票/平均，降低方差。
    \item \textbf{Boosting (GBDT, XGBoost, LightGBM)}：多棵树串行训练，每棵修正前一棵的残差，降低偏差。
    \item 集成学习是当前表格数据竞赛中的\textbf{绝对统治者}。
\end{itemize}

\section{关键代码片段}
\begin{lstlisting}
from sklearn.tree import DecisionTreeClassifier, export_text
from sklearn.datasets import load_iris
from sklearn.model_selection import train_test_split

# 加载数据
X, y = load_iris(return_X_y=True)
X_train, X_test, y_train, y_test = train_test_split(X, y, test_size=0.3)

# 构建 CART 决策树
tree = DecisionTreeClassifier(
    criterion='gini',       # 基尼指数
    max_depth=3,             # 预剪枝: 最大深度
    min_samples_leaf=5       # 预剪枝: 叶子最小样本数
)
tree.fit(X_train, y_train)
print(f"Accuracy: {tree.score(X_test, y_test):.4f}")

# 打印树的文本表示
print(export_text(tree, feature_names=load_iris().feature_names))
\end{lstlisting}

\section{深度面试题}

\textbf{问：为什么 Random Forest 通常比单棵决策树好？}\\
答：单棵决策树是高方差模型（不稳定，数据微小变化会导致完全不同的树）。Random Forest 通过 1. Bagging（有放回采样）降低方差；2. 特征随机选择进一步降低树之间的相关性。根据方差分解公式：$\text{Var}(\bar{X}) = \frac{\rho \sigma^2 + (1-\rho)\sigma^2/n}{1}$，当树之间相关性 $\rho$ 降低时，集成的方差大幅下降。

\textbf{问：决策树处理缺失值的策略是什么？}\\
答：C4.5 的方案是：在计算信息增益时，对缺失值样本按比例分配到各个子节点（权重按非缺失样本的分布比例确定）。CART 使用代理变量（Surrogate Split）：选择一个与原始最优特征最相似的替代特征来处理缺失值。

\section{参考文献与拓展}
\begin{enumerate}
    \item \textbf{ID3}: Quinlan, J. R. (1986). \textit{Induction of decision trees}. Machine Learning.
    \item \textbf{C4.5}: Quinlan, J. R. (1993). \textit{C4.5: Programs for Machine Learning}. Morgan Kaufmann.
    \item \textbf{CART}: Breiman, L., et al. (1984). \textit{Classification and Regression Trees}. Wadsworth.
    \item \textbf{教材}: 周志华《机器学习》第 4 章（决策树）。
\end{enumerate}

\end{document}
